\documentclass[10pt]{extarticle}
\usepackage[a4paper,margin=1.5cm]{geometry}
\usepackage[T1]{fontenc}
\usepackage{amsmath,amssymb}
\usepackage{graphicx}
\usepackage[export]{adjustbox}
\usepackage{multicol}
\usepackage{enumitem}
\usepackage{booktabs,array,tabularx}
\setlength{\columnsep}{0.7cm}
\setlength{\parindent}{0pt}
\setlength{\parskip}{3pt}
\pagestyle{empty}
\makeatletter
\newcommand{\qfigdef}[4]{\@namedef{qf@file@#1}{#2}\@namedef{qf@w@#1}{#3}\@namedef{qf@wide@#1}{#4}}
\newcommand{\qf@check}[1]{\@ifundefined{qf@file@#1}{\PackageError{score2}{Unknown figure id #1}{}}{}}
\newcommand{\qf@img}[4][]{\includegraphics[width=#2\dimexpr\@nameuse{qf@w@#4}pt\relax,max width=#3,max height=0.45\textheight,#1]{\@nameuse{qf@file@#4}}}
\newif\ifqf@full
\newif\ifqf@boxed
\newsavebox{\qf@box}
\newcommand{\qf@open}{\global\qf@fulltrue\par\global\qf@boxedtrue
  \setbox\qf@box\vbox\bgroup\hsize\textwidth\linewidth\textwidth}
\newcommand{\qf@emit}[1]{\par\egroup\global\qf@boxedfalse
  \ifdim\dimexpr\ht\qf@box+\dp\qf@box\relax<#1\textheight\noindent\usebox{\qf@box}\par\else\unvbox\qf@box\fi}
\newcommand{\qfullsplit}{\ifqf@boxed\qf@emit{0.85}\fi}
\newcommand{\qf@close}{\ifqf@boxed\qf@emit{0.45}\fi\global\qf@fullfalse}
\newcommand{\qfullbegin}{\end{multicols}\qf@open}
\newcommand{\qfullend}{\qf@close\begin{multicols}{2}}
\newcommand{\qfullstart}{\qf@open}
\newcommand{\qfulldone}{\qf@close}
\newcommand{\qfig}[2][1]{\qf@check{#2}%
  \ifqf@full\qf@blk{#1}{#2}\else
  \ifnum\@nameuse{qf@wide@#2}=1 \expandafter\@firstoftwo\else\expandafter\@secondoftwo\fi
  {\end{multicols}\begin{center}\qf@img{#1}{\textwidth}{#2}\end{center}\begin{multicols}{2}}%
  {\qf@blk{#1}{#2}}\fi}
\newcommand{\qf@blk}[2]{\par\smallskip{\centering\qf@img{#1}{\linewidth}{#2}\par}\smallskip}
\newcommand{\qfigi}[2][1]{\qf@check{#2}\qf@img[valign=c]{#1}{\linewidth}{#2}}
\makeatother
\newcommand{\papersec}[1]{\par\bigskip{\centering\Large\bfseries PAPER-#1\par}\medskip}
\newcommand{\partsec}[2]{\par\medskip\noindent{\large\bfseries PART-#1: #2}\par\smallskip}
\newcommand{\sect}[1]{\par\smallskip\noindent\textbf{#1}\par}
\newcommand{\qn}[1]{\par\smallskip\noindent\textbf{#1.}~}
\newcommand{\paraq}[2]{\par\smallskip\noindent\textbf{Paragraph for Questions #1 and #2}\par}
\newenvironment{qoptions}{\begin{itemize}[leftmargin=1.8em,itemsep=1pt,topsep=2pt,label={}]}{\end{itemize}}
\newcommand{\opt}[2]{\item[(#1)] #2}
% figure registry (generated by score2_tools.py finalize)
\qfigdef{p1_phy_s1_q01_1}{Score-II - Test 31 (29 Mar 2024)-img/p1_phy_s1_q01_1.png}{124.2}{0}
\qfigdef{p1_phy_s1_q04_1}{Score-II - Test 31 (29 Mar 2024)-img/p1_phy_s1_q04_1.png}{68.6}{0}
\qfigdef{p1_phy_s1_q05_1}{Score-II - Test 31 (29 Mar 2024)-img/p1_phy_s1_q05_1.png}{138.9}{0}
\qfigdef{p1_phy_s1_q08_1}{Score-II - Test 31 (29 Mar 2024)-img/p1_phy_s1_q08_1.png}{95.1}{0}
\qfigdef{p1_phy_s1_q12_1}{Score-II - Test 31 (29 Mar 2024)-img/p1_phy_s1_q12_1.png}{116.1}{0}
\qfigdef{p1_phy_s1_q13_1}{Score-II - Test 31 (29 Mar 2024)-img/p1_phy_s1_q13_1.png}{93.0}{0}
\qfigdef{p1_phy_s1_q16_1}{Score-II - Test 31 (29 Mar 2024)-img/p1_phy_s1_q16_1.png}{68.8}{0}
\qfigdef{p1_phy_s1_q16_2}{Score-II - Test 31 (29 Mar 2024)-img/p1_phy_s1_q16_2.png}{66.3}{0}
\qfigdef{p1_phy_s1_q16_3}{Score-II - Test 31 (29 Mar 2024)-img/p1_phy_s1_q16_3.png}{66.3}{0}
\qfigdef{p1_phy_s1_q16_4}{Score-II - Test 31 (29 Mar 2024)-img/p1_phy_s1_q16_4.png}{66.3}{0}
\qfigdef{p1_phy_s1_q16_5}{Score-II - Test 31 (29 Mar 2024)-img/p1_phy_s1_q16_5.png}{66.3}{0}
\qfigdef{p1_phy_s1_q17_1}{Score-II - Test 31 (29 Mar 2024)-img/p1_phy_s1_q17_1.png}{129.4}{0}
\qfigdef{p1_phy_s1_q17_2}{Score-II - Test 31 (29 Mar 2024)-img/p1_phy_s1_q17_2.png}{116.3}{0}
\qfigdef{p1_phy_s1_q18_1}{Score-II - Test 31 (29 Mar 2024)-img/p1_phy_s1_q18_1.png}{82.0}{0}
\qfigdef{p1_phy_s1_q19_1}{Score-II - Test 31 (29 Mar 2024)-img/p1_phy_s1_q19_1.png}{108.2}{0}
\qfigdef{p1_phy_s1_q20_1}{Score-II - Test 31 (29 Mar 2024)-img/p1_phy_s1_q20_1.png}{100.6}{0}
\qfigdef{p1_phy_s2_q04_1}{Score-II - Test 31 (29 Mar 2024)-img/p1_phy_s2_q04_1.png}{81.1}{0}
\qfigdef{p1_phy_s2_q05_1}{Score-II - Test 31 (29 Mar 2024)-img/p1_phy_s2_q05_1.png}{115.9}{0}
\qfigdef{p1_phy_s2_q06_1}{Score-II - Test 31 (29 Mar 2024)-img/p1_phy_s2_q06_1.png}{78.9}{0}
\qfigdef{p1_chem_s1_q08_1}{Score-II - Test 31 (29 Mar 2024)-img/p1_chem_s1_q08_1.png}{145.8}{0}
\qfigdef{p1_chem_s1_q09_1}{Score-II - Test 31 (29 Mar 2024)-img/p1_chem_s1_q09_1.png}{21.6}{0}
\qfigdef{p1_chem_s1_q11_1}{Score-II - Test 31 (29 Mar 2024)-img/p1_chem_s1_q11_1.png}{114.4}{0}
\qfigdef{p1_chem_s1_q11_3}{Score-II - Test 31 (29 Mar 2024)-img/p1_chem_s1_q11_3.png}{47.2}{0}
\qfigdef{p1_chem_s1_q11_4}{Score-II - Test 31 (29 Mar 2024)-img/p1_chem_s1_q11_4.png}{35.0}{0}
\qfigdef{p1_chem_s1_q11_5}{Score-II - Test 31 (29 Mar 2024)-img/p1_chem_s1_q11_5.png}{39.4}{0}
\qfigdef{p1_chem_s1_q11_6}{Score-II - Test 31 (29 Mar 2024)-img/p1_chem_s1_q11_6.png}{21.1}{0}
\qfigdef{p1_chem_s1_q12_1}{Score-II - Test 31 (29 Mar 2024)-img/p1_chem_s1_q12_1.png}{156.8}{0}
\qfigdef{p1_chem_s1_q13_1}{Score-II - Test 31 (29 Mar 2024)-img/p1_chem_s1_q13_1.png}{155.0}{0}
\qfigdef{p1_chem_s1_q13_2}{Score-II - Test 31 (29 Mar 2024)-img/p1_chem_s1_q13_2.png}{46.6}{0}
\qfigdef{p1_chem_s1_q13_3}{Score-II - Test 31 (29 Mar 2024)-img/p1_chem_s1_q13_3.png}{45.1}{0}
\qfigdef{p1_chem_s1_q13_4}{Score-II - Test 31 (29 Mar 2024)-img/p1_chem_s1_q13_4.png}{49.4}{0}
\qfigdef{p1_chem_s1_q13_5}{Score-II - Test 31 (29 Mar 2024)-img/p1_chem_s1_q13_5.png}{44.8}{0}
\qfigdef{p1_chem_s2_q04_2}{Score-II - Test 31 (29 Mar 2024)-img/p1_chem_s2_q04_2.png}{50.2}{0}
\qfigdef{p1_chem_s2_q04_1}{Score-II - Test 31 (29 Mar 2024)-img/p1_chem_s2_q04_1.png}{43.2}{0}
\qfigdef{p1_chem_s2_q04_4}{Score-II - Test 31 (29 Mar 2024)-img/p1_chem_s2_q04_4.png}{49.4}{0}
\qfigdef{p1_chem_s2_q04_3}{Score-II - Test 31 (29 Mar 2024)-img/p1_chem_s2_q04_3.png}{35.8}{0}
\qfigdef{p1_chem_s2_q04_5}{Score-II - Test 31 (29 Mar 2024)-img/p1_chem_s2_q04_5.png}{141.3}{0}
\qfigdef{p1_chem_s2_q05_1}{Score-II - Test 31 (29 Mar 2024)-img/p1_chem_s2_q05_1.png}{141.3}{0}
\qfigdef{p1_chem_s2_q06_1}{Score-II - Test 31 (29 Mar 2024)-img/p1_chem_s2_q06_1.png}{41.3}{0}
\qfigdef{p1_chem_s2_q06_2}{Score-II - Test 31 (29 Mar 2024)-img/p1_chem_s2_q06_2.png}{70.6}{0}
\qfigdef{p1_chem_s2_q06_3}{Score-II - Test 31 (29 Mar 2024)-img/p1_chem_s2_q06_3.png}{86.2}{0}
\qfigdef{p1_chem_s2_q06_4}{Score-II - Test 31 (29 Mar 2024)-img/p1_chem_s2_q06_4.png}{35.5}{0}
\qfigdef{p1_chem_s2_q07_1}{Score-II - Test 31 (29 Mar 2024)-img/p1_chem_s2_q07_1.png}{158.5}{0}
\qfigdef{p1_math_s1_q15_1}{Score-II - Test 31 (29 Mar 2024)-img/p1_math_s1_q15_1.png}{139.1}{0}
\makeatletter\@beginparpenalty=10000\makeatother
\begin{document}
\begin{center}{\Large\bfseries JEE (Main+Advanced) Score-II --- Test 31 (29 Mar 2024)}\end{center}
\begin{multicols}{2}
\raggedcolumns
\partsec{1}{PHYSICS}
\sect{Section-I}

\qn{1} An object of mass $m=1$ kg is attached to a platform of mass $M=4$ kg with spring of spring constant $k=400$ N/m. There is no friction between the object and the platform, and the coefficient of static friction between the platform and the ground is $\mu=0.1$. The object is placed at its equilibrium position, and then given a horizontal velocity $v$. The value of $v$ for which platform is just about to slip on the ground, is :-
\qfig{p1_phy_s1_q01_1}
\begin{qoptions}
\opt{A}{0.4 m/s}
\opt{B}{0.25 m/s}
\opt{C}{0.5 m/s}
\opt{D}{0.1 m/s}
\end{qoptions}

\qn{2} If $\eta$ represents the coefficient of viscosity and $T$ the surface tension, then the dimension of $\dfrac T\eta$ is same as that of :
\begin{qoptions}
\opt{A}{length}
\opt{B}{mass}
\opt{C}{time}
\opt{D}{speed}
\end{qoptions}

\qn{3} \textbf{Assertion:} If a gas container in motion is suddenly stopped, the temperature of the gas rises.\\
\textbf{Reason:} The kinetic energy of ordered mechanical motion is converted into the kinetic energy of random motion of gas molecules.
\begin{qoptions}
\opt{A}{Assertion is correct, Reason is correct; Reason is a correct explanation for Assertion.}
\opt{B}{Assertion is correct, Reason is correct, Reason is not a correct explanation for Assertion}
\opt{C}{Assertion is correct, Reason is incorrect}
\opt{D}{Assertion is incorrect, Reason is correct}
\end{qoptions}

\qn{4} A long rope is designed in such a way that $v$ is independent of $Z$ (distance from point of suspension) as shown in the figure. For such a rope $\mu$ (linear mass density of rope) cannot be constant. General equation for $\mu(Z)$ is :- [use $\mu_0$ as mass per unit length of the rope at point $Z=0$]
\qfig{p1_phy_s1_q04_1}
\begin{qoptions}
\opt{A}{$\mu=4\mu_0e^{-(g/v^2)z}$}
\opt{B}{$\mu=\mu_0e^{-(g/v^2)z}$}
\opt{C}{$\mu=4\mu_0e^{(g/v^2)z}$}
\opt{D}{$\mu=\mu_0e^{(g/v^4)z}$}
\end{qoptions}

\qn{5} Shown in figure a certain quantities of kerosene and water is filled in hemispherical bowl. The specific gravity of kerosene is 0.8. The force exerted by the water on kerosene is approximately :-\\
(Take atmospheric pressure $=10^5$ pa)
\qfig{p1_phy_s1_q05_1}
\begin{qoptions}
\opt{A}{$2\times10^7$ N}
\opt{B}{$1.85\times10^7$ N}
\opt{C}{$3\times10^7$ N}
\opt{D}{$2.5\times10^7$ N}
\end{qoptions}

\qn{6} A student performs an experiment to determine the young's modulus of a wire, exactly 2 m long, by Searle method. In a particular reading, the student measures the extension in the length of the wire to be 0.8 mm with an uncertainty of $\pm0.05$ mm at a load of exactly 1.0 kg. The student also measures the diameter of the wire to be 0.4 mm with an uncertainty of $\pm0.01$ mm. Take : $g=9.8\ \mathrm{m/s^2}$ (exact). The maximum permissible error in measurement of Young's modulus from the experiment is :-
\begin{qoptions}
\opt{A}{5.6 \%}
\opt{B}{11.25\%}
\opt{C}{2.5\%}
\opt{D}{6\%}
\end{qoptions}

\qn{7} Consider an ideal gas confined in an isolated closed chamber. As the gas undergoes an adiabatic expansion, the average time of collision between molecules increases as $T^q$, where $T$ is the temperature of the gas. The value of $q$ is :- $\left(\gamma=\dfrac{C_p}{C_v}\right)$
\begin{qoptions}
\opt{A}{$\dfrac{1-\gamma}{2}$}
\opt{B}{$\dfrac{1-2\gamma}{2}$}
\opt{C}{$\dfrac{1+\gamma}{2(1-\gamma)}$}
\opt{D}{$\dfrac{1-\gamma}{2(1+\gamma)}$}
\end{qoptions}

\qn{8} Two particles A and B are dropped simultaneously from rest both $h$ m above the ground. Particle A falls on the ground, while the particle B hits the smooth inclined surface in its mid path as shown in figure. The ratio of time of fall of the projectile B to the time to fall of A to ground is $\left(e=\dfrac13\right)$ :-
\qfig{p1_phy_s1_q08_1}
\begin{qoptions}
\opt{A}{$\sqrt2$}
\opt{B}{$\dfrac1{\sqrt2}$}
\opt{C}{$\sqrt3$}
\opt{D}{3}
\end{qoptions}

\qn{9} Light of wavelength 488 nm is produced by an argon laser which is used in the photoelectric effect. When light from this spectral line is incident on the emitter, the stopping (cut-off) potential of photoelectrons is 0.38 V. Find the work function of the material from which the emitter is made.
\begin{qoptions}
\opt{A}{2.17 eV}
\opt{B}{3.16 eV}
\opt{C}{4.01 eV}
\opt{D}{2.92 eV}
\end{qoptions}

\qn{10} The curve of binding energy per nucleon as a functions of atomic mass number has a sharp peak for helium .This implies that helium
\begin{qoptions}
\opt{A}{can easily be broke up}
\opt{B}{is very stable}
\opt{C}{can be used as fissionable material}
\opt{D}{is radioactive}
\end{qoptions}

\qn{11} An air plane, with 20 m wingspread is flying at 250 ms$^{-1}$ parallel to the earth's surface at a place where the horizontal component of earth's magnetic field is $2\times10^{-5}$ T and angle of dip is 60$^\circ$. The magnitude of the induced emf between the tips of the wings is:
\begin{qoptions}
\opt{A}{$\dfrac1{10}$ V}
\opt{B}{$\dfrac{\sqrt2}{10}$ V}
\opt{C}{$\dfrac{\sqrt3}{10}$ V}
\opt{D}{$\dfrac15$ V}
\end{qoptions}

\qn{12} For given current-time graph, find ratio $\dfrac{i_{\text{avg}}}{i_{\text{rms}}}$.
\qfig{p1_phy_s1_q12_1}
\begin{qoptions}
\opt{A}{tan45$^\circ$}
\opt{B}{sin30$^\circ$}
\opt{C}{cos30$^\circ$}
\opt{D}{sec30$^\circ$}
\end{qoptions}

\qn{13} A conducting wire is in the shape of a regular hexagon, which is inscribed inside an imaginary circle of radius $R$. as shown. A current `I' flows through the wire. The magnitude of the magnetic field at the center of the circle is :-
\qfig{p1_phy_s1_q13_1}
\begin{qoptions}
\opt{A}{$\dfrac{3\mu_0I}{2\pi R}$}
\opt{B}{$\dfrac{\sqrt3\mu_0I}{2\pi R}$}
\opt{C}{$\dfrac{\mu_0I}{2\sqrt3\pi R}$}
\opt{D}{$\dfrac{\sqrt3\mu_0I}{\pi R}$}
\end{qoptions}

\qn{14} Two polaroids are placed in the path of unpolarized beam of intensity $I_0$ such that no light is emitted from the second polaroid. If a third polaroid whose polarization axis makes an angle $\theta$ with the polarization axis of first polaroid, is placed between these polaroids, then the intensity of light emerging from the last polaroid will be :-
\begin{qoptions}
\opt{A}{$\left(\dfrac{I_0}8\right)\sin^22\theta$}
\opt{B}{$\left(\dfrac{I_0}4\right)\sin^22\theta$}
\opt{C}{$\left(\dfrac{I_0}2\right)\cos^2\theta$}
\opt{D}{$I_0\cos^4\theta$}
\end{qoptions}

\qn{15} \textbf{Statement$-$1 :} When a width of one of the slits of Young's double slit experiment is four times that of the other, the brighter fringes are nine times brighter than the dark fringes.\\
and\\
\textbf{Statement$-$2 :} The amplitude of the wave is proportional to the width of the slit.
\begin{qoptions}
\opt{A}{Statement$-$1 is True, Statement$-$2 is True ; Statement$-$2 is a correct explanation for Statement$-$1.}
\opt{B}{Statement$-$1 is True, Statement$-$2 is True ; Statement$-$2 is not a correct explanation for Statement$-$1.}
\opt{C}{Statement$-$1 is True, Statement$-$2 is False.}
\opt{D}{Statement$-$1 is False, Statement$-$2 is True.}
\end{qoptions}

\qn{16} Wavefronts incident on an interface between the media are shown in the figure. The refracted wavefront will be as shown in :-
\qfig{p1_phy_s1_q16_1}
\begin{qoptions}
\opt{A}{\qfigi{p1_phy_s1_q16_2}}
\opt{B}{\qfigi{p1_phy_s1_q16_3}}
\opt{C}{\qfigi{p1_phy_s1_q16_4}}
\opt{D}{\qfigi{p1_phy_s1_q16_5}}
\end{qoptions}

\qn{17} Here signals A, B, C are shown, then the logic gate is :-
\qfig{p1_phy_s1_q17_1}
\qfig{p1_phy_s1_q17_2}
\begin{qoptions}
\opt{A}{AND Gate}
\opt{B}{OR Gate}
\opt{C}{NOR Gate}
\opt{D}{NAND Gate}
\end{qoptions}

\qn{18} Consider a spherical drop of radius $R$. Surface tension of the liquid is $S$. Force of surface tension on shaded sub hemisphere due to remaining drop is $\dfrac{S\pi R}2$. The value of angle $\theta$ (angle subtended by the sub hemisphere at centre of drop) is :-
\qfig{p1_phy_s1_q18_1}
\begin{qoptions}
\opt{A}{60$^\circ$}
\opt{B}{120$^\circ$}
\opt{C}{30$^\circ$}
\opt{D}{90$^\circ$}
\end{qoptions}

\qn{19} Figure above shows a closed Gaussian surface in the shape of a cube of edge length 3.0 m. There exists an electric field given by $[(2.0x+4.0)\hat i+8.0\,\hat j+3.0\,\hat k]$ N/C, where $x$ is in meters, in the region in which it lies. The net charge in coulombs enclosed by the cube is equal to :-
\qfig{p1_phy_s1_q19_1}
\begin{qoptions}
\opt{A}{$-54\,\epsilon_0$}
\opt{B}{$6\,\epsilon_0$}
\opt{C}{$-6\,\epsilon_0$}
\opt{D}{$54\,\epsilon_0$}
\end{qoptions}

\qn{20} A stationary point object (O) is at a distance $h$ from water surface in which a plane mirror is placed at a depth $d$. An observer sees the image of object in mirror (as shown). If the mirror starts moving upward with constant speed $v$ then speed of image formed by plane mirror as observed by observer will be :-
\qfig{p1_phy_s1_q20_1}
\begin{qoptions}
\opt{A}{$v$ upwards}
\opt{B}{$\dfrac32v$ upwards}
\opt{C}{$2v$ upwards}
\opt{D}{$v$ downwards}
\end{qoptions}

\sect{Section-II}

\qn{1} The time period of simple pendulum is `$T$'. When the length increases by 10 cm, its period is $T_1$. When the length is decreased by 10 cm, it's period is $T_2$. Then the relation between $T$, $T_1$ and $T_2$ is given by $\alpha T^2=T_1^\beta+T_2^\gamma$ then $\alpha+\beta+\gamma$ is :-

\qn{2} The magnifying power of the telescope is found to be 9 and the separation between the lenses is 20 cm for relaxed eye. Then magnitude of difference of focal lengths (in cm) of component lenses is :-

\qn{3} 64 tuning forks are arranged in increasing order of frequency such that each fork produces 4 beats per second with next one. If the frequency of the last fork is double of the first fork, the frequency (in Hz) of 16th fork is :-

\qn{4} A cubical block of side $a$ and density $\rho$ slides inside the groove of a fixed incline plane with constant velocity $v$. There is a thin film of viscous fluid of thickness $t$ between surface of plane and the block. Friction coefficient between cube and side walls of groove is $\mu$. Viscosity of thin film is given by $\dfrac{N\rho agt}{10V}\left(1-\sqrt3\,\mu\right)$. Find $N$.
\qfig{p1_phy_s2_q04_1}

\qn{5} In given circuit, power (in W) of zener diode for which bulb brights at rated value (Nearly) is :-
\qfig{p1_phy_s2_q05_1}

\qn{6} For the given track, optimum speed is 30m/s then maximum speed is $20\sqrt x$ m/s for turn, then the value $x$ is :-
\qfig{p1_phy_s2_q06_1}

\qn{7} Two particles execute S.H.M about same mean position. The position of particle is given by $x_1=10\sin\left(\pi t+\dfrac\pi6\right)$cm, $x_2=10\cos(\pi t)$cm. Then maximum separation (in cm) between them during motion will be :-

\qn{8} An observer moves away from the linear source of sound. The loudness level decreases by 20 dB, when the distance of observer from source is increases by $10^x$ times, then the value of $x$ is :-

\qn{9} An electron in hydrogen atom revolves around its nucleus with speed $12\times10^6$ m/s in an orbit of radius 0.52\AA. The magnetic field (in T) produced at the centre of the hydrogen atom is :-

\qn{10} A potentiometer wire has length 4m and resistance 8$\Omega$. The resistance (in $\Omega$) that must be connected in series with the wire and an accumulator of e.m.f 2V, so as to get a potential gradient 1mV per cm on the wire is :-

\partsec{2}{CHEMISTRY}
\sect{Section-I}

\qn{1} Identify the \textbf{INCORRECT} option regarding oxide of nitrogen is :
\begin{qoptions}
\opt{A}{In N$_2$O$_5$ , all N$-$O bond lengths are not equal.}
\opt{B}{In N$_2$O$_4$ , all N$-$O bond lengths are equal.}
\opt{C}{O$\widehat{\text{N}}$O bond angle in NO$_2$ is nearly same when it will dimerise.}
\opt{D}{In N$_2$O , all bond length N$-$N \& N$-$O bond lengths are equal.}
\end{qoptions}

\qn{2} Select the \textbf{CORRECT} order according to give property.
\begin{qoptions}
\opt{A}{Lewis acidic nature : BCl$_3$ $<$ AlCl$_3$ $<$ GaCl$_3$}
\opt{B}{Response for magnet : CaSO$_4$ $<$ FeSO$_4$ $<$ MnSO$_4$}
\opt{C}{Reducing property : AsH$_3$ $<$ PH$_3$ $<$ NH$_3$}
\opt{D}{Solubility : AgF $<$ AgCl $<$ AgBr}
\end{qoptions}

\qn{3} Assume group elements according to data in p-block.
\par\nobreak
\begin{center}
\begin{tabular}{|p{1.7cm}|c|c|c|c|c|}\hline
\footnotesize$\Delta_{\text{eg}}$H (kJ mol$^{-1}$) & \footnotesize$-141$ & \footnotesize$-200$ & \footnotesize$-195$ & \footnotesize$-190$ & \footnotesize$-174$\\\hline
\footnotesize Group elements (A to E top to bottom) & A & B & C & D & E\\\hline
\end{tabular}
\end{center}
Consider the following statements.\\
(i) All elements belongs to 15th group\\
(ii) Hydrides of all elements shows hydrogen bonding\\
(iii) Element `B' and `C' can form acidic oxide.\\
(iv) Element `A' can not show allotropy.\\
Choose correct option :
\begin{qoptions}
\opt{A}{(i), (ii) (iii)}
\opt{B}{(i), (iv)}
\opt{C}{(i), (iii), (iv)}
\opt{D}{only (iii)}
\end{qoptions}

\qn{4} Which of the following bond length (in pm) is correct for B$_2$ , C$_2$ , N$_2$ respectively?
\begin{qoptions}
\opt{A}{159, 131, 110}
\opt{B}{110, 143, 131}
\opt{C}{110, 131, 143}
\opt{D}{159, 110, 131}
\end{qoptions}

\qn{5} \textbf{Statement-I :} Ruby is colored due to 0.5 $-$ 1\% of Al$^{3+}$ ions\\
\textbf{Statement-II :} Ruby is colored due to electronic d$-$d transition in octahedral chromium(III) complex.
\begin{qoptions}
\opt{A}{Both Statement-I and Statement-II are correct.}
\opt{B}{Statement-I is correct but Statement-II is incorrect.}
\opt{C}{Both Statement-I and Statement-II are incorrect.}
\opt{D}{Statement-I is incorrect but Statement-II is correct.}
\end{qoptions}

\qn{6} Which of the following statement is correct?
\begin{qoptions}
\opt{A}{FeO$_4^{-2}$ ferrates are formed in acidic medium but they readily decompose to Fe$_2$O$_3$ and O$_2$}
\opt{B}{The E$^\Theta$ value for the Mn$^{+3}$/Mn$^{+2}$ $\Rightarrow$ +1.57 is much more positive than Fe$^{+3}$/Fe$^{+2}$ $\Rightarrow$ +0.77 is mainly responsible for third ionisation energy of Mn}
\opt{C}{Fe$^{+2}$ stronger reducing agent than Cr$^{+2}$}
\opt{D}{In all the 3d series metals, Zn has smallest radius.}
\end{qoptions}

\qn{7} Which of the following element has half filled `f' subshell in its MO$_2$ type oxide?
\begin{qoptions}
\opt{A}{Pr (Z = 59)}
\opt{B}{Nd (Z = 60)}
\opt{C}{Dy (Z = 66)}
\opt{D}{Tb (Z = 65)}
\end{qoptions}

\qn{8} Which of the following is/are non aromatic in nature ?
\qfig{p1_chem_s1_q08_1}
\begin{qoptions}
\opt{A}{Only A}
\opt{B}{A and C}
\opt{C}{A, C and D}
\opt{D}{A, B, C and D}
\end{qoptions}

\qn{9} Compound formed D is an isomers of which compound.
\begin{center}\qfigi{p1_chem_s1_q09_1}\ $\xrightarrow[0^\text{o}\text{C}]{\text{NaNO}_2/\text{HCl}}$ A $\xrightarrow{\text{H}_2\text{O}/\Delta}$ B $\xrightarrow[\text{(ii) Dil.HCl}]{\text{(i) CHCl}_3/\text{KOH}}$ C\\[2pt] $\xrightarrow[\text{(ii) H}^{+}/\text{H}_2\text{O}]{\text{(i) 1 Eq.Cl}_2/h\nu}$ D\end{center}
\begin{qoptions}
\opt{A}{Vanillin}
\opt{B}{Cinnamaldehyde}
\opt{C}{Camphor}
\opt{D}{Terpineol}
\end{qoptions}

\qn{10} Identify incorrect statement from the following :
\begin{qoptions}
\opt{A}{All monosaccharides are reducing sugar}
\opt{B}{Glucose don't give schiff's test.}
\opt{C}{In sucrose $\alpha$, $\beta$-glycosidic linkage is present.}
\opt{D}{Amylopectin is insoluble in water and having unbranched chain polymer.}
\end{qoptions}

\qn{11} Compound D gives lucas reagent test Identify compound E.
\qfig{p1_chem_s1_q11_1}
\begin{center}$\xrightarrow[\text{BaSO}_4]{\text{H}_2/\text{Pd}}$ C $\xrightarrow[\text{NaOH}]{\text{Conc.}}$ E + D\end{center}
\begin{qoptions}
\opt{A}{\qfigi{p1_chem_s1_q11_3}}
\opt{B}{\qfigi{p1_chem_s1_q11_4}}
\opt{C}{\qfigi{p1_chem_s1_q11_5}}
\opt{D}{\qfigi{p1_chem_s1_q11_6}}
\end{qoptions}

\qn{12} Match the column
\qfig{p1_chem_s1_q12_1}
\begin{qoptions}
\opt{A}{A-I, B-IV, C-V, D-III}
\opt{B}{A-II, B-V, C-IV, D-I}
\opt{C}{A-II, B-I, C-IV, D-III}
\opt{D}{A-III, B-I, C-IV, D-II}
\end{qoptions}

\qn{13} \qfig{p1_chem_s1_q13_1}
\begin{qoptions}
\opt{A}{\qfigi{p1_chem_s1_q13_2}}
\opt{B}{\qfigi{p1_chem_s1_q13_3}}
\opt{C}{\qfigi{p1_chem_s1_q13_4}}
\opt{D}{\qfigi{p1_chem_s1_q13_5}}
\end{qoptions}

\qn{14} The density of a solution prepared by dissolving 120 g of urea (molar mass = 60 g) in 1000 g of water is 1.15 g/mL. The molarity of this solution is :
\begin{qoptions}
\opt{A}{0.50 M}
\opt{B}{2.05 M}
\opt{C}{1.02 M}
\opt{D}{1.78 M}
\end{qoptions}

\qn{15} A photon was absorbed by a hydrogen atom in its ground state, and the electron was prompted to the fifth orbit. When the excited atom returned to its ground state. In this process, how many maximum spectral lines could be obtained $-$
\begin{qoptions}
\opt{A}{1}
\opt{B}{2}
\opt{C}{5}
\opt{D}{10}
\end{qoptions}

\qn{16} The vapour pressure of the solution of two liquids A($P^0=80$ mm Hg) and B($P^0=120$ mm Hg) is found to be 100 mm Hg when $X_A=0.4$. The result shows that
\begin{qoptions}
\opt{A}{Solution exhibits ideal behaviour}
\opt{B}{Solution shows positive deviation from ideal behaviour}
\opt{C}{Solution shows negative deviation from ideal behaviour}
\opt{D}{Attractive intermolecular forces between A$-$A in pure liquid A and B$-$B in pure liquid B are stronger than those between A$-$B when mixed in solution.}
\end{qoptions}

\qn{17} At 30$^\circ$C, the half-life for the decomposition of AB$_2$ is 200 sec and is independent of the initial concentration of AB$_2$. The time required (in sec) for 80\% of the AB$_2$ to decompose is (Given: log 2 = 0.30; log 3 = 0.48)
\begin{qoptions}
\opt{A}{200 sec}
\opt{B}{323 sec}
\opt{C}{467 sec}
\opt{D}{532 sec}
\end{qoptions}

\qn{18} Calculate $\Delta$H when 2 moles of solid benzoic acid undergo complete combustion at 300K if :\\
C$_6$H$_5$COOH(s) + $\dfrac{15}2$ O$_2$(g) $\longrightarrow$ 7CO$_2$(g) + 3H$_2$O($l$)\\
$\Delta$U$_{\text{reaction}}$ = $-$ 750 kJ/mole\\
\relax[R = 8 Joule/ mole-K]
\begin{qoptions}
\opt{A}{$-$ 751.2 kJ}
\opt{B}{$-$ 752.4 kJ}
\opt{C}{$-$ 1501.2 kJ}
\opt{D}{$-$ 1502.4 kJ}
\end{qoptions}

\qn{19} Three solution of strong electrolytes 25 ml of 0.1 M HX, 25 ml of 0.1 M H$_2$Y and 50 ml of 0.1N Z(OH)$_2$ are mixed. What is the pOH of the resulting solution at 25$^\circ$C ? [Given log 5 = 0.7]
\begin{qoptions}
\opt{A}{1.6}
\opt{B}{7}
\opt{C}{11.6}
\opt{D}{12.4}
\end{qoptions}

\qn{20} \textbf{Assertion :} The SRP of three metallic ions A$^{+}$, B$^{2+}$, C$^{3+}$ are $-0.3$, $-0.5$, 0.8 volt respectively, so oxidising power of ions is C$^{3+}$ $>$ A$^{+}$ $>$ B$^{2+}$.\\
\textbf{Reason :} Higher the SRP, higher the oxidising power.
\begin{qoptions}
\opt{A}{Both assertion and reason are true and the reason is a correct explanation of assertion.}
\opt{B}{Both assertion and reason are true but reason is not a correct explanation of assertion.}
\opt{C}{Assertion is true but the reason is false.}
\opt{D}{Assertion is false but the reason is true.}
\end{qoptions}

\sect{Section-II}

\qn{1} The octahedral complex $[\mathrm{Ti(H_2O)_6}]^{+3}$ has a single d-electron. To excite this electron from the ground state t$_{2g}$ orbital to the e$_g$ orbital, this complex absorb light from 450 to 600 nm. The maximum absorbance corresponds to $\Delta_{\text{oct.}}$ and occurs at 499 nm. the value of $\Delta_{\text{oct.}}$ is `a' $\times10^{-19}$ Joule\\
Find the value of `a' in nearest integer.\\
(hc = 1240 ev/nm)

\qn{2} Consider the following order according to indicate property.\\
(i) B$-$F bond length : BF$_3$ $<$ BF$_2$OH\\
(ii) Dipole moment : CHCl$_3$ $<$ CH$_3$Cl\\
(iii) Acidic strength : SO$_3$ $<$ SeO$_3$\\
(iv) Number of d-orbital , participate in hybridisation : BrF$_5$ $<$ PF$_5$\\
Find the total number of correct orders?

\qn{3} Consider the following cations :\\
Pb$^{+2}$ , Cu$^{+2}$ , As$^{+3}$ , A$\ell^{+3}$ , Fe$^{+3}$ , Co$^{+2}$ , Ni$^{+2}$ , Mn$^{+2}$ , Zn$^{+2}$\\
Find the total Number of cations which produce black ppt. with (NH$_4$)$_2$S

\qn{4} How many of the following compounds don't give racemic mixture with CH$_3$MgBr :
\par\nobreak
\begin{center}\small
\begin{tabular}{|c|c|c|c|}\hline
(A) & \qfigi{p1_chem_s2_q04_2} & (B) & \qfigi{p1_chem_s2_q04_1}\\\hline
(C) & \qfigi{p1_chem_s2_q04_4} & (D) & \qfigi{p1_chem_s2_q04_3}\\\hline
\multicolumn{4}{|c|}{\qfigi{p1_chem_s2_q04_5}}\\\hline
\end{tabular}
\end{center}

\qn{5} Total number of sp$^2$ hybridised orbital in A :\\
\qfig{p1_chem_s2_q05_1}

\qn{6} The total number of amines which can't be synthesized by Gabriel phthalimide is :\\
(I) CH$_3-$NH$-$CH$_3$\\
\qfigi{p1_chem_s2_q06_1}\\
\qfigi{p1_chem_s2_q06_2}\\
\qfigi{p1_chem_s2_q06_3}\\
\qfigi{p1_chem_s2_q06_4}

\qn{7} How many of the following compounds don't give racemic mixture with CH$_3$MgBr :
\qfig{p1_chem_s2_q07_1}

\qn{8} A(g) + 2B(s) $\rightleftharpoons$ 2C(g).\\
Initially 2 mol A(g), 4 mole of B(s) and 1 mole of a inert gas are present in a closed container. After equilibrium has established total pressure of container becomes 9 atm. If A(g) consume 50\% at equilibrium then calculate Kp (in atm) for above reaction -

\qn{9} Consider an isothermal reversible compression of one mole of an ideal gas in which the pressure of the system is increased from 5 atm to 30 atm at 300 K. Calculate the entropy change of the surrounding (in J K$^{-1}$).\\
\relax[Given: R = $\dfrac{25}3$ J mol$^{-1}$ K$^{-1}$, ln 6 = 1.8]

\qn{10} What volume (in mL) of a 1.0 M HCl solution should be added to in aqueous solution of 1.0 L 0.001 M HCl in order to decrease its pH by one unit ?

\partsec{3}{MATHEMATICS}
\sect{Section-I}

\qn{1} Let $Z=x+iy$ where $x,y\in R$ where $i=\sqrt{-1}$ then the minimum value of the expression $|z+2|+|z-2-i|+|z-3|+|z-1+2i|$
\begin{qoptions}
\opt{A}{$\sqrt5+\sqrt{10}$}
\opt{B}{$10-\sqrt5$}
\opt{C}{$\sqrt5+10$}
\opt{D}{$5+\sqrt{10}$}
\end{qoptions}

\qn{2} Let A be square matrix of order $4\times4$, using entries only 0 \& 1 then the probability that each and every row and column must contain exactly one `1'.
\begin{qoptions}
\opt{A}{$\dfrac3{2^{13}}$}
\opt{B}{$\left(\dfrac5{2^{13}}\right)$}
\opt{C}{$\left(\dfrac3{2^{12}}\right)$}
\opt{D}{None of these}
\end{qoptions}

\qn{3} \textbf{Assertion : -} Number of ways in which 75600 can be resolved as a product of two divisors is 60.\\
\textbf{Reason : -} If total divisors of $N$ is `$t$' then number of ways in which $N$ can be resolved as a product of two divisors is $\dfrac t2$.
\begin{qoptions}
\opt{A}{Both assertion and reason are true and reason is correct explanation of assertion.}
\opt{B}{Both assertion and reason are true and reason is not correct explanation of assertion.}
\opt{C}{Assertion is true but reason is false.}
\opt{D}{Assertion is false.}
\end{qoptions}

\qn{4} If $f(x)=\cos x\cos(x+2)-\cos^2(x+1)$ then the value of $f(9)$ is -
\begin{qoptions}
\opt{A}{$\cos2-\cos^21$}
\opt{B}{$\cos4-\cos^22$}
\opt{C}{$\cos8-\cos^210$}
\opt{D}{None of these}
\end{qoptions}

\qn{5} The area of a sector of an ellipse $\dfrac{x^2}{16}+\dfrac{y^2}9=1$ between the major axis and a radius vector from the focus whose eccentric angle $\theta$ is
\begin{qoptions}
\opt{A}{$\dfrac34\left(4\theta-\sqrt7\sin\theta\right)$}
\opt{B}{$\dfrac34\left(2\theta-\sqrt7\sin\theta\right)$}
\opt{C}{$\dfrac34\left(2\theta-\dfrac{\sqrt7}2\sin\theta\right)$}
\opt{D}{$\dfrac34\left(8\theta-2\sqrt7\sin\theta\right)$}
\end{qoptions}

\qn{6} $f(x)=\sin^4x+\cos^4x+bx+c$. The values of $b$ so that $f(x)$ is strictly monotonic
\begin{qoptions}
\opt{A}{$b\leq-1$ or $b\geq1$}
\opt{B}{$-1\leq b\leq-1$}
\opt{C}{$b\geq0$ or $b\leq-1$}
\opt{D}{$b\leq0$ or $b\geq1$}
\end{qoptions}

\qn{7} The value of $\displaystyle\sum_{r=0}^n\dfrac{{}^{n+1}C_r}{{}^{n+1}C_r+{}^{n+1}C_{r+1}}$ is
\begin{qoptions}
\opt{A}{$n+2$}
\opt{B}{$\dfrac{n+3}2$}
\opt{C}{$\dfrac{n+1}2$}
\opt{D}{$\dfrac{n+2}2$}
\end{qoptions}

\qn{8} Let $a_n$ and $b_n$ be $n$th terms of two arithmetic progressions where $n\in N$ and their sum of first $n$-terms be $S_a(n)$ and $S_b(n)$ respectively.\\
Given $\dfrac{S_a(n)}{S_b(n)}=\dfrac{5n+9}{2n+8}$ and $a_2=4$ then the value of $\dfrac{a_1}{b_1}$ is
\begin{qoptions}
\opt{A}{3/5}
\opt{B}{4/5}
\opt{C}{6/5}
\opt{D}{7/5}
\end{qoptions}

\qn{9} Let $f(x)=\sin x+\displaystyle\int_{-\pi/2}^{\pi/2}(\sin x+t\cos x)f(t)\,dt$ then match the entries of List-I with List-II and mark the correct option -
\par\nobreak
\begin{center}\small
\begin{tabular}{|c|p{3.2cm}|c|c|}\hline
 & \multicolumn{1}{c|}{\textbf{List-I}} & & \textbf{List-II}\\\hline
(A) & If range of $f(x)$ is $[\alpha,\beta]$ then $\alpha+\beta$ is & (I) & $\dfrac23$\\\hline
(B) & $f\left(\dfrac{3\pi}2\right)$ is & (II) & $\dfrac1{\sqrt2}$\\\hline
(C) & $f(\pi)$ is & (III) & $\dfrac13$\\\hline
(D) & $f\left(\dfrac{5\pi}4\right)$ is & (IV) & 0\\\hline
\end{tabular}
\end{center}
\begin{qoptions}
\opt{A}{(A) $-$ III, (B) $-$ I, (C) $-$ IV, (D) $-$ II}
\opt{B}{(A) $-$ III, (B) $-$ II, (C) $-$ I, (D) $-$ IV}
\opt{C}{(A) $-$ I, (B) $-$ III, (C) $-$ IV, (D) $-$ II}
\opt{D}{(A) $-$ IV, (B) $-$ III, (C) $-$ I, (D) $-$ II}
\end{qoptions}

\qn{10} The locus of the mid-points of the perpendiculars drawn from points on the line $x=\lambda y$, to the line $x=y$, is $5x-7y=0$, then the value of $\lambda$ is -
\begin{qoptions}
\opt{A}{2}
\opt{B}{3}
\opt{C}{4}
\opt{D}{6}
\end{qoptions}

\qn{11} If the ellipse $\dfrac{x^2}{16}+\dfrac{y^2}7=1$ and the hyperbola $\dfrac{x^2}{144}-\dfrac{y^2}\alpha=\dfrac1{25}$ are orthogonal, then find the value of $4e_1$. (where $e_1$ is the eccentricity of the hyperbola)
\begin{qoptions}
\opt{A}{10}
\opt{B}{5}
\opt{C}{8}
\opt{D}{6}
\end{qoptions}

\qn{12} In how many ways two subsets B and C can be formed from the set A = \{1, 2, 3, 4, 5, 6\} independently such that B $\cup$ C = A.
\begin{qoptions}
\opt{A}{243}
\opt{B}{729}
\opt{C}{64}
\opt{D}{120}
\end{qoptions}

\qn{13} The angle between the straight lines, whose direction cosines $\ell$, $m$, $n$ are given by the equations $2\ell+2m-n=0$ and $mn+n\ell+\ell m=0$ is $\dfrac\pi k$, then the sum of all the divisors of $6k^2$
\begin{qoptions}
\opt{A}{120}
\opt{B}{30}
\opt{C}{60}
\opt{D}{90}
\end{qoptions}

\qn{14} Let $F(x)$ be a function defined as\\
$F(x)=\begin{cases}\max\left\{t^3-12t+5\right\} & ;\ x\leqslant3,\ t\leqslant x\\ [x+2]+4 & ;\ x\in(3,5]\\ 3x+5 & ;\ x>5\end{cases}$\\
where $[\ ]$ is greatest integer function. Let $m$ and $n$ be number of points of discontinuity and non-differentiability of $F(x)$ respectively then $(m+n)$ lies in
\begin{qoptions}
\opt{A}{$(-\infty,-2]$}
\opt{B}{$[-2,2]$}
\opt{C}{$[2,5]$}
\opt{D}{$[5,\infty)$}
\end{qoptions}

\qn{15} Let $f(x)$ be a continous function in $[0,7]$ and graph of $f'(x)$ is given below
\qfig{p1_math_s1_q15_1}
$S_1$ : $f(x)$ has three points of extrema\\
$S_2$ : $f(x)$ has 4 points of inflection
\begin{qoptions}
\opt{A}{Both statement-I and statement-II are correct.}
\opt{B}{Both statement-I and statement-II are incorrect.}
\opt{C}{Statement-I is incorrect and statement-II is correct.}
\opt{D}{Statement-I is correct and statement-II is incorrect.}
\end{qoptions}

\qn{16} If $f(x)=2^x+x^3+1$ and $g(x)=f^{-1}(x)$ then value of $\dfrac{g'(2)}{g''(2)}$ lies in
\begin{qoptions}
\opt{A}{Set of irrationals numbers}
\opt{B}{Set of prime numbers}
\opt{C}{$[-10,10]$}
\opt{D}{None of these}
\end{qoptions}

\qn{17} Let $I(x)=\displaystyle\int\dfrac{x^2\left(x\cos\mathrm{ec}^2x-\cot x\right)}{(1-x\cot x)^2}dx$. If $I\left(\dfrac\pi2\right)=0$ then $I\left(\dfrac\pi4\right)$ is equal to
\begin{qoptions}
\opt{A}{$\dfrac{\pi^2}{4(\pi+4)}-2\,\ell n\left(\dfrac{4+\pi}{4\sqrt2}\right)+\dfrac{\pi^2}8$}
\opt{B}{$\dfrac{\pi^2}{4(\pi-4)}+2\,\ell n\left(\dfrac{4-\pi}{4\sqrt2}\right)$}
\opt{C}{$\dfrac{\pi^2}{4(\pi+4)}-2\,\ell n\left(\dfrac{4+\pi}{4\sqrt2}\right)+\dfrac{\pi^2}4$}
\opt{D}{$\dfrac{\pi^2}{4(\pi-4)}+2\,\ell n\left(\dfrac{4-\pi}{4\sqrt2}\right)+\dfrac{\pi^2}4$}
\end{qoptions}

\qn{18} If $x^2-4x+\log_{\frac14}\left(a^2\right)=0$ does not have two distinct real roots, then minimum value of $a$ is
\begin{qoptions}
\opt{A}{1/16}
\opt{B}{$-1/4$}
\opt{C}{$-1/16$}
\opt{D}{1/4}
\end{qoptions}

\qn{19} Let $x,y,z$ and $u$ be positive integers such that $\log_xy=\dfrac32$, $\log_zu=\dfrac54$, If $x-z=9$ then $y-u=$
\begin{qoptions}
\opt{A}{3}
\opt{B}{21}
\opt{C}{93}
\opt{D}{None}
\end{qoptions}

\qn{20} $f(x)=x^3+ax^2+bx-6$ is a 3 degree polynomial. If $f(x)$ has its local minima and local maxima at $\alpha+\sqrt\beta$ and $\alpha-\sqrt\beta$ respectively such that $\left(\alpha-\sqrt\beta\right)>0$. If $\alpha=2$ and all 3 real roots are positive integers, then value of $a+b+\beta^{-1}$ is
\begin{qoptions}
\opt{A}{2}
\opt{B}{4}
\opt{C}{8}
\opt{D}{6}
\end{qoptions}

\sect{Section-II}

\qn{1} If the value of the series\\ $\left\{\dfrac1{\sin1^\circ\sin2^\circ}+\dfrac1{\sin2^\circ\sin3^\circ}+\dfrac1{\sin3^\circ\sin4^\circ}+\cdots\right.$\\ $\left.\cdots+\dfrac1{\sin89^\circ\sin90^\circ}\right\}$ is $\dfrac{\cos(a^\circ)}{\sin^2(b^\circ)}$ then the value of $a+b$ is \underline{\hspace{2em}}.

\qn{2} Let $f:[-1,1]\to\left[\dfrac{5\pi}2,\dfrac{7\pi}2\right]$ where $f(x)=\sin^{-1}x$ such that $\sin^{-1}(\sin5)=a\pi-b$ then $(2a-b)$ is \underline{\hspace{2em}}

\qn{3} Let A be a set of first 10 odd natural numbers and B be a set of first 10 even natural numbers then number of onto functions $f:A\to B$ such that $f(1)>f(3)>f(5)>f(7)<f(9)<f(11)<f(13)$ is/are N then the value of N/144

\qn{4} Let A and B are square matrices of third order such that $A^TB^2=A$ ($A^T$ is transpose of A)\\
$S_1$ : $\left|A-A^T\right|=0$\\
$S_2$ : $\left|A^2+\left(A^T\right)^2\right|=0$\\
$S_3$ : $\left|B^2\left(B^T\right)^2\right|=0$\\
$S_4$ : $\left|B^2+\left(B^T\right)^2\right|=0$\\
Number of statement is/are true -

\qn{5} The smallest possible natural number `$p$' for which the equation $x^2-px+1258=0$ has integral roots ?

\qn{6} Let $(x+\alpha)^2=-k(y+\beta)$ be the equation of the mirror image of the parabola $y^2=16x$ with respect to the line $x+y+4=0$, then the value of $\dfrac{k}{2(\alpha+\beta)}$ is

\qn{7} If $x,y,z$ be three distinct real numbers, and $x\neq1$, $y\neq1$, $z\neq1$. If the vectors $\vec p=x\hat i+\hat j+\hat k$, $\vec q=\hat i+y\hat j+\hat k$, $\vec r=\hat i+\hat j+z\hat k$ are such that $(\vec p\times\vec q)\cdot\vec r=0$, then $\dfrac1{1-x}+\dfrac1{1-y}+\dfrac1{1-z}$ is equal to

\qn{8} Let $a\in(0,\infty)$. Let $P:R\to R$ and $Q:(a,\infty)\to R$ be the functions defined by $P(x)=\log_{16}\left(\dfrac{ex}2\right)$ and $Q(x)=\dfrac1{2e}\dfrac{\ell n(x-a)}{\ell n\left(e^x-e^a\right)}$ then value of $100\left|\displaystyle\lim_{x\to a^+}P(Q(x))\right|$ is

\qn{9} Let $I(x)=\displaystyle\int\dfrac{\sec^2x-\beta}{\sin^\alpha x}dx$ and $g(x)=I(x)-(\alpha-\beta)\displaystyle\int\sin^{-\alpha}x\,dx$. The value of $g\left(\dfrac\pi4\right)$ at $\alpha=20$.

\qn{10} $S$ represents a curve whose length of the normal at any point P is equal to the radius vector of the point P, passes through (5, 3) is $x^2\pm y^2=c$, then sum of all possible values of $c$ is

\end{multicols}
\end{document}
